\documentclass[11pt]{article}

\usepackage[margin=1in]{geometry}
\usepackage{amsmath,amssymb,amsthm}
\usepackage{graphicx}
\usepackage[hidelinks]{hyperref}

\newtheorem{theorem}{Theorem}
\newtheorem{lemma}{Lemma}
\newtheorem{remark}{Remark}

\newcommand{\R}{\mathbb{R}}
\newcommand{\norm}[1]{\left\lVert #1 \right\rVert}
\newcommand{\weak}{\xrightarrow{w}}

\title{Positive Multilinear Forms on Non-Atomic \texorpdfstring{$L_p$}{Lp}}
\author{Automated Math Discovery Smoke Test}
\date{June 8, 2026}

\begin{document}
\maketitle

\section*{Source and Open Question}

This packet concerns Question 3 from:

\begin{quote}
Luis A. Garcia, Jose Lucas P. Luiz, Vinicius C. C. Miranda,
\emph{Norm attainment for multilinear operators and polynomials on
Banach Spaces}\\
\emph{and Banach lattices}, arXiv:2605.12117.
\end{quote}

The source location is \texttt{arxiv.tex:1007--1071}, final section
``Open Questions'', Question 3. The question appears on page 28 of the
source paper:

\begin{center}
\includegraphics[width=\textwidth]{figures/open_problem_crop.png}
\end{center}

The paper defines an $n$-linear operator $A:X_1\times\cdots\times X_n\to Y$
to be weakly sequentially continuous when
\[
  A(x_{1,k},\ldots,x_{n,k})\to A(x_1,\ldots,x_n)
\]
whenever $x_{j,k}\weak x_j$ in $X_j$ for every $j=1,\ldots,n$.

\begin{remark}[Exponent convention]
The rendered question writes $L_{p_i}[0,1]$ and assumes
$\sum_i 1/p_i<1$. We read this in the surrounding finite-exponent
convention of the paper, namely $1<p_i<\infty$. The inequality already
excludes $p_i=1$. A human reviewer should confirm whether the authors
intended to include the endpoint $p_i=\infty$; the proof below resolves
the finite-exponent formulation.
\end{remark}

\section*{Claim}

\begin{theorem}
Let $n\ge 2$ and let $1<p_1,\ldots,p_n<\infty$ satisfy
\[
  \sum_{i=1}^n \frac{1}{p_i}<1.
\]
Then there is a positive continuous $n$-linear form on
\[
  L_{p_1}[0,1]\times\cdots\times L_{p_n}[0,1]
\]
which is not weakly sequentially continuous. In particular, Question 3 has
a negative answer under the finite-exponent interpretation.
\end{theorem}

\begin{proof}
Define
\[
  A(f_1,\ldots,f_n)=\int_0^1 f_1(t)\cdots f_n(t)\,dt.
\]
Since $\sum_i 1/p_i<1$, choose $s\in(1,\infty]$ such that
\[
  \frac1s=1-\sum_{i=1}^n \frac1{p_i}.
\]
Holder's inequality, with the remaining factor filled by the constant
function $1\in L_s[0,1]$, gives
\[
  |A(f_1,\ldots,f_n)|
  \le \prod_{i=1}^n \norm{f_i}_{p_i}.
\]
Thus $A$ is a continuous $n$-linear form. It is positive in the real
Banach-lattice sense, because $f_i\ge 0$ for all $i$ implies
$f_1\cdots f_n\ge 0$ almost everywhere.

Let $(r_k)$ be the Rademacher functions on $[0,1]$. For every
$1<p<\infty$, one has $r_k\weak 0$ in $L_p[0,1]$. Indeed, if $q$ is the
conjugate exponent and $h\in L_q[0,1]$, approximate $h$ in $L_q$ by a
dyadic step function $h_m$. For all sufficiently large $k$,
$\int_0^1 h_m r_k=0$, and Holder's inequality controls
$\int_0^1 (h-h_m)r_k$.

Now set
\[
  x_{1,k}=r_k,\qquad x_{2,k}=r_k,\qquad x_{j,k}=1\quad (3\le j\le n).
\]
Then
\[
  x_{1,k}\weak 0,\qquad x_{2,k}\weak 0,\qquad x_{j,k}\weak 1
  \quad (3\le j\le n),
\]
but
\[
  A(r_k,r_k,1,\ldots,1)=\int_0^1 r_k(t)^2\,dt=1
\]
for every $k$, while
\[
  A(0,0,1,\ldots,1)=0.
\]
Therefore $A$ is not weakly sequentially continuous.
\end{proof}

\section*{Verification Notes}

No computational verification is needed. The proof reduces to three
standard facts: Holder's inequality, positivity of the integral form on
real $L_p$ lattices, and weak nullity of the Rademacher sequence in finite
$L_p[0,1]$.

The only convention-sensitive point is whether Question 3 intended the
endpoint $p_i=\infty$. The packet should be read as a full negative answer
for the finite-exponent version used in the surrounding paper.

\end{document}
